% Options for packages loaded elsewhere
\PassOptionsToPackage{unicode}{hyperref}
\PassOptionsToPackage{hyphens}{url}
\PassOptionsToPackage{dvipsnames,svgnames,x11names}{xcolor}
%
\documentclass[
  letterpaper,
  DIV=11,
  numbers=noendperiod]{scrartcl}

\usepackage{amsmath,amssymb}
\usepackage{iftex}
\ifPDFTeX
  \usepackage[T1]{fontenc}
  \usepackage[utf8]{inputenc}
  \usepackage{textcomp} % provide euro and other symbols
\else % if luatex or xetex
  \usepackage{unicode-math}
  \defaultfontfeatures{Scale=MatchLowercase}
  \defaultfontfeatures[\rmfamily]{Ligatures=TeX,Scale=1}
\fi
\usepackage{lmodern}
\ifPDFTeX\else  
    % xetex/luatex font selection
\fi
% Use upquote if available, for straight quotes in verbatim environments
\IfFileExists{upquote.sty}{\usepackage{upquote}}{}
\IfFileExists{microtype.sty}{% use microtype if available
  \usepackage[]{microtype}
  \UseMicrotypeSet[protrusion]{basicmath} % disable protrusion for tt fonts
}{}
\makeatletter
\@ifundefined{KOMAClassName}{% if non-KOMA class
  \IfFileExists{parskip.sty}{%
    \usepackage{parskip}
  }{% else
    \setlength{\parindent}{0pt}
    \setlength{\parskip}{6pt plus 2pt minus 1pt}}
}{% if KOMA class
  \KOMAoptions{parskip=half}}
\makeatother
\usepackage{xcolor}
\setlength{\emergencystretch}{3em} % prevent overfull lines
\setcounter{secnumdepth}{-\maxdimen} % remove section numbering
% Make \paragraph and \subparagraph free-standing
\makeatletter
\ifx\paragraph\undefined\else
  \let\oldparagraph\paragraph
  \renewcommand{\paragraph}{
    \@ifstar
      \xxxParagraphStar
      \xxxParagraphNoStar
  }
  \newcommand{\xxxParagraphStar}[1]{\oldparagraph*{#1}\mbox{}}
  \newcommand{\xxxParagraphNoStar}[1]{\oldparagraph{#1}\mbox{}}
\fi
\ifx\subparagraph\undefined\else
  \let\oldsubparagraph\subparagraph
  \renewcommand{\subparagraph}{
    \@ifstar
      \xxxSubParagraphStar
      \xxxSubParagraphNoStar
  }
  \newcommand{\xxxSubParagraphStar}[1]{\oldsubparagraph*{#1}\mbox{}}
  \newcommand{\xxxSubParagraphNoStar}[1]{\oldsubparagraph{#1}\mbox{}}
\fi
\makeatother


\providecommand{\tightlist}{%
  \setlength{\itemsep}{0pt}\setlength{\parskip}{0pt}}\usepackage{longtable,booktabs,array}
\usepackage{calc} % for calculating minipage widths
% Correct order of tables after \paragraph or \subparagraph
\usepackage{etoolbox}
\makeatletter
\patchcmd\longtable{\par}{\if@noskipsec\mbox{}\fi\par}{}{}
\makeatother
% Allow footnotes in longtable head/foot
\IfFileExists{footnotehyper.sty}{\usepackage{footnotehyper}}{\usepackage{footnote}}
\makesavenoteenv{longtable}
\usepackage{graphicx}
\makeatletter
\newsavebox\pandoc@box
\newcommand*\pandocbounded[1]{% scales image to fit in text height/width
  \sbox\pandoc@box{#1}%
  \Gscale@div\@tempa{\textheight}{\dimexpr\ht\pandoc@box+\dp\pandoc@box\relax}%
  \Gscale@div\@tempb{\linewidth}{\wd\pandoc@box}%
  \ifdim\@tempb\p@<\@tempa\p@\let\@tempa\@tempb\fi% select the smaller of both
  \ifdim\@tempa\p@<\p@\scalebox{\@tempa}{\usebox\pandoc@box}%
  \else\usebox{\pandoc@box}%
  \fi%
}
% Set default figure placement to htbp
\def\fps@figure{htbp}
\makeatother

\KOMAoption{captions}{tableheading}
\makeatletter
\@ifpackageloaded{caption}{}{\usepackage{caption}}
\AtBeginDocument{%
\ifdefined\contentsname
  \renewcommand*\contentsname{Table of contents}
\else
  \newcommand\contentsname{Table of contents}
\fi
\ifdefined\listfigurename
  \renewcommand*\listfigurename{List of Figures}
\else
  \newcommand\listfigurename{List of Figures}
\fi
\ifdefined\listtablename
  \renewcommand*\listtablename{List of Tables}
\else
  \newcommand\listtablename{List of Tables}
\fi
\ifdefined\figurename
  \renewcommand*\figurename{Figure}
\else
  \newcommand\figurename{Figure}
\fi
\ifdefined\tablename
  \renewcommand*\tablename{Table}
\else
  \newcommand\tablename{Table}
\fi
}
\@ifpackageloaded{float}{}{\usepackage{float}}
\floatstyle{ruled}
\@ifundefined{c@chapter}{\newfloat{codelisting}{h}{lop}}{\newfloat{codelisting}{h}{lop}[chapter]}
\floatname{codelisting}{Listing}
\newcommand*\listoflistings{\listof{codelisting}{List of Listings}}
\makeatother
\makeatletter
\makeatother
\makeatletter
\@ifpackageloaded{caption}{}{\usepackage{caption}}
\@ifpackageloaded{subcaption}{}{\usepackage{subcaption}}
\makeatother

\usepackage{bookmark}

\IfFileExists{xurl.sty}{\usepackage{xurl}}{} % add URL line breaks if available
\urlstyle{same} % disable monospaced font for URLs
\hypersetup{
  colorlinks=true,
  linkcolor={blue},
  filecolor={Maroon},
  citecolor={Blue},
  urlcolor={Blue},
  pdfcreator={LaTeX via pandoc}}


\author{}
\date{}

\begin{document}


\section{1.14}\label{section}

\(P(AB)\) is the probability of the interesection of the event sets
\(A\) and \(B\).

Thus the intersection of the two sets is at most the size of the smaller
set. Thus \(P(AB)\) is at most \(P(A)\)

To find the minimum probability, its trying to minimize the intersection
of sets \(A\) and \(B\). In otherwords, maximumizing the intersection of
\(A\) and \(B^c\). \(P(A)=0.4\) and \(P(B^c)=1-P(B)=1-0.7=0.3\) The
maximum of this intersection is \(0.3\) because \(P(B^c)\) is the
smaller one, thus the remaining \(0.1\) probability must be in the
intersection of \(A\) and \((B^{c})^{c}=B\). Thus the minimum of
\(P(AB)\) is \(0.1\)

Therefore \(0.1 \leq P(AB) \leq 0.4\)

\section{1.18}\label{section-1}

\(\Omega = \{3,4,5\}\)

\begin{longtable}[]{@{}clccc@{}}
\toprule\noalign{}
\(k\) & \textbar{} & \(3\) & \(4\) & \(5\) \\
\midrule\noalign{}
\endhead
\bottomrule\noalign{}
\endlastfoot
\(p_X(k) = P(X == k)\) & \textbar{} & \(\frac3{16}\) & \(\frac12\) &
\(\frac{5}{16}\) \\
\end{longtable}

\section{1.26}\label{section-2}

\(|\Omega | = \binom{15}{4}  =1365\)

\(|\text{2 men 2 women}| = \binom{10}{2} * \binom{5}{2}= 45*10=450\)

\(P(\text{2 men 2 women}) = \frac{|\text{2m2w}|}{|\Omega|} =\frac{450}{1365}=\frac{30}{91}\)

\section{1.28}\label{section-3}

\subsection{a}\label{a}

\(A\) is w/o replacement

\(P(A)= \frac{n^2-n + m^2-m}{(n+m)(n+m-1)}\)

\subsection{b}\label{b}

\(B\) is w/ replacement

\(P(B)=\frac{n^2 + m^2}{(n+m)^2}\)

\subsection{c}\label{c}

Intuitively, when you do not replace the ball, the probability you get
the same color again is lower because there is less of the same color
now. Therefore \(P(A)\) is smaller than \(P(B)\).

\section{1.33}\label{section-4}

\(|\Omega| = 6^5\)

\(|FH| = 6*5\binom 53=300\)

\(P(FH) = \frac{300}{6^5}=\frac{25}{648}\)

\section{1.41}\label{section-5}

\(A_i=\{\text{player } i \text{ wins no games}\}\)

\(P(A_i) = \frac23 ^4= \frac{16}{81}\)

\(A_i \cap A_j=\{\text{players }i,j \text{ win no games}\}\)

\(P(A_i\cap A_j) =\frac{1}{3^4}= \frac{1}{81}\)

\(A_1 \cap A_2 \cap A_3=\{\text{no one wins any games}\}\)

\(P(A_1\cap A_2 \cap A_3) = 0\) because at least one person must win
each round.

\(P(A_1 \cup A_2 \cup A_3) = \sum_{k=1}^3(-1)^{k+1} \sum_{1\leq i_1 < ... < i_k \leq n} P(A_{i1}\cap ... \cap A_{ik}) = 3P(A_i) + 3P(A_i\cap A_j) + P(A_1\cap A_2 \cap A_3)=\frac{16}{27} + \frac{1}{27} + 0= \frac{17}{27}\)

\section{1.44 (tentative)}\label{tentative}

\subsection{a}\label{a-1}

\(X,Y\in \{1,2,3,4,5,6\}\)

\subsection{b and c}\label{b-and-c}

\begin{longtable}[]{@{}
  >{\centering\arraybackslash}p{(\linewidth - 14\tabcolsep) * \real{0.1600}}
  >{\centering\arraybackslash}p{(\linewidth - 14\tabcolsep) * \real{0.1200}}
  >{\centering\arraybackslash}p{(\linewidth - 14\tabcolsep) * \real{0.1200}}
  >{\centering\arraybackslash}p{(\linewidth - 14\tabcolsep) * \real{0.1200}}
  >{\centering\arraybackslash}p{(\linewidth - 14\tabcolsep) * \real{0.1200}}
  >{\centering\arraybackslash}p{(\linewidth - 14\tabcolsep) * \real{0.1200}}
  >{\centering\arraybackslash}p{(\linewidth - 14\tabcolsep) * \real{0.1200}}
  >{\centering\arraybackslash}p{(\linewidth - 14\tabcolsep) * \real{0.1200}}@{}}
\toprule\noalign{}
\begin{minipage}[b]{\linewidth}\centering
\(k\)
\end{minipage} & \begin{minipage}[b]{\linewidth}\centering
\textbar{}
\end{minipage} & \begin{minipage}[b]{\linewidth}\centering
\(1\)
\end{minipage} & \begin{minipage}[b]{\linewidth}\centering
\(2\)
\end{minipage} & \begin{minipage}[b]{\linewidth}\centering
\(3\)
\end{minipage} & \begin{minipage}[b]{\linewidth}\centering
\(4\)
\end{minipage} & \begin{minipage}[b]{\linewidth}\centering
\(5\)
\end{minipage} & \begin{minipage}[b]{\linewidth}\centering
\(6\)
\end{minipage} \\
\midrule\noalign{}
\endhead
\bottomrule\noalign{}
\endlastfoot
\(P(X=k)\) & \textbar{} & \(\frac{1}{36}\) & \(\frac{1}{12}\) &
\(\frac{5}{36}\) & \(\frac{7}{36}\) & \(\frac{9}{36}\) &
\(\frac{11}{36}\) \\
\(P(X\leq k)\) & \textbar{} & \(\frac{1}{36}\) & \(\frac{1}{9}\) &
\(\frac{1}{4}\) & \(\frac{4}{9}\) & \(\frac{25}{36}\) & \(1\) \\
\(P(Y =  k)\) & \textbar{} & \(\frac{11}{36}\) & \(\frac{9}{36}\) &
\(\frac{7}{36}\) & \(\frac{5}{36}\) & \(\frac{3}{36}\) &
\(\frac{1}{36}\) \\
\(P(Y\leq k)\) & \textbar{} & \(\frac{11}{36}\) & \(\frac{5}{9}\) &
\(\frac{3}{4}\) & \(\frac{8}{9}\) & \(\frac{35}{36}\) & \(1\) \\
\end{longtable}




\end{document}
